% a01_H0rel part 4: Dirichlet-face positivity, the face family, the positivity (3.26) at the weight
\subsection{The Dirichlet-face positivity and the face family}\label{app:a01:face}

\begin{lemma}[Dirichlet-face positivity]\label{lem:a01:FPC}
In unit-$j$ units let $y'\in\Lone$ be a face block, $B(y')$ its unit cube, and $\Delta^{(j)}$ the block form of the cube at any gauge field $A$. Then for every block field $\phi$ and $\theta\in(0,1)$,
\[
\ip\phi{\Delta^{(j)}\phi}\ \ge\ (1-\theta)\ip\phi{\Delta^{(j)}\phi}+\theta\,\gamma_D\,|\phi(y')|^2,\qquad\gamma_D=\frac{a_j}{1+a_j\Wt^2}.
\]
The same holds for the composite multiscale form of \cite[(3.23)--(3.25)]{Ba2} restricted to the fine sites of $B^j(y')$, in the form
\[
\sum_{\text{bonds of }\overline{B}}\eps^{d}|D_A\phi_0|^2+\sum_{i<j}a\,\ell_{i+1}^{d-2}\sum_{B^{i+1}(y'')\subset B}|\phi_{i+1}(y'')-Q(A)\phi_i(y'')|^2\ \ge\ (\Wt+a_j^{-1})^{-1}\ell_j^{d-2}|\phi_j(y')|^2
\]
(physical units; the first sum over the bonds inside $B=B^j(y')$ and the face bonds from $B$ to the killed layer), for every gauge field.
\end{lemma}
\begin{proof}
By the minimiser identity, $\ip\phi{\Delta\phi}=\min_\Phi F(\Phi)$, $F=\ip\Phi{(-\Delta^D_A+V_j)\Phi}+a_j\norm{\phi-Q\Phi}^2$. In $\theta F$ keep only the gradient terms of bonds inside $B(y')$, the Dirichlet terms of the face bonds of $B(y')$ (each $\vth^{d-2}|\Phi(x)|^2$) and the block term of $y'$; all dropped terms are $\ge0$. With $g\defeq|\Phi|$, unitary transport gives $|D_A\Phi(b)|\ge|g(b_+)-g(b_-)|$, and the face terms are $\vth^{d-2}g(x)^2$. Discrete Friedrichs inequality on the unit cube with the killed layer behind one face: on each line normal to that face, $g(x)=\sum_{s\le m}[g(s)-g(s-1)]$ with $g(0)\defeq0$ at the killed site and $m\le\vth^{-1}$, so $g(x)^2\le\vth^{-1}\sum_{\text{line}}|\partial g|^2$; summing $\vth^dg^2$ over the line and over all lines, $\sum_{B(y')}\vth^dg^2\le E(g)$, the kept energy. Then $|(Q\Phi)(y')|\le\sum\vth^dw|\Phi|\le\Wt(\sum\vth^dg^2)^{1/2}\le\Wt E^{1/2}$. So $\theta F\ge\theta\min_{s\ge0}[\Wt^{-2}s^2+a_j(|\phi(y')|-s)_+^2]=\theta\gamma_D|\phi(y')|^2$ (springs in series), and minimising over $\Phi$ gives the first claim. Composite form: with $s_i\defeq(\ell_i^d\sum_B|\phi_i|^2)^{1/2}$, Friedrichs along the normal lines from the killed layer gives $s_0^2\le\ell_j^2E$ with $E$ the first sum; Cauchy--Schwarz with the weights gives $\norm{Q_1\phi_0}\le\Wt^{1/2}s_0$ at the weighted step and $\norm{Q\phi_i}\le s_i$ at the uniform steps $i\ge1$; the reverse triangle inequality gives $s_j\le\Wt^{1/2}s_0+\sum_it_i$, $t_i\defeq\norm{\phi_{i+1}-Q\phi_i}$, and Cauchy--Schwarz with the weights $(\Wt\ell_j^2,\ell_{i+1}^2/a)$ gives $s_j^2\le(\Wt\ell_j^2+\ell_j^2/a_j)(E+\sum_ia\ell_{i+1}^{-2}t_i^2)$, using $\sum_{i<j}\ell_{i+1}^2/a=\ell_j^2/a_j$.
\end{proof}

This is the extra term $\gamma_0\sum_{x\in\bar\Omega^{(k)}}|\phi(x)|^2$ of \cite[Prop.~5.4, (5.26)]{K2}, valid for every gauge field (King's printed proof \cite[(5.28)--(5.29)]{K2} expands around a constant field and needs $A$ regular). King's constant \cite[(5.38)--(5.40)]{K2} is consistent with it.

\begin{definition}[the face family]\label{def:a01:facefam}
At every scale $0\le i\le K_f$ of the upper bound's large-field decomposition, $i=0$ included, the \emph{face family} is the set of events $\{|\phi_i(y')|>\sigma_i\}$ at $y'\in\Lone^{(i)}\cap\Lambda_7^{(i-1)\prime}$ (with $\Lambda_7^{(-1)\prime}\defeq$ everything), where $\Lone^{(i)}$ is the set of $i$-blocks containing a site of $F_1$ (for $i\ge1$ the face blocks of both faces in every direction; for $i=0$ the first live layer $F_1$ itself, where $\phi_0$ is the fine field), and $\sigma_i\defeq(2/(\theta_2\gamma_D^{(i)}))^{1/2}p_i\le C_8p_i$ with $\gamma_D^{(0)}\defeq1$. At the stopping scale the same test is the indicator $\lambda_j$ of \cref{app:a01:weights} (it is \cite[(5.41)]{K2} with $C=\sigma_j/p_j$).
\end{definition}

At $i=0$ the payment is the $\theta_2$-share of the face-bond term itself: $\ip{\phi_0}{(-\Delta^D_A+V)\phi_0}\ge\sum_{\text{face bonds }b\ni x}|D_A\phi_0(b)|^2=n_F(x)|\phi_0(x)|^2$ (unit-$0$ units; no link on a face bond), with $n_F(x)\in\{1,2,3\}$ the number of face bonds at $x$. The scale-$0$ member must be stated on the first live layer (not as the set of boundary blocks, which at $i=0$ would be the killed plane on one face), because the inherited cap of \cref{lem:a01:IC} at $k=1$ needs a cap on the first live layer of both faces.

\begin{lemma}[payment, and the corrected near-face estimates]\label{lem:a01:F541}
(a) A violation of the face test at $y'$ at scale $i$ is paid by the factor $e^{-p_i^2/4}$ from the $\theta_2$-share of the scale-$i$ form (\cref{lem:a01:FPC}), and in the composite form of \cite[Prop.~3.1]{Ba2} by its composite clause; violating blocks of different scales are disjoint as fine sets by Ba{\l}aban's region rules, so the payments add. (b) For $x\in\Lambda_4^{(i)}$ with $\dist(x,\DD)\le r(\ell_i)$, a coordinate block path from $x$ to the face layer of length $\le r(\ell_i)$ lies in $\Lambda_3^{(i)}\subset\Lambda_0^{(i)}$ (since $\dist(\Lambda_4,\Lambda_3^c)>r$), so on the complement of the face family and of the gradient family $|\phi_i(x)|\le\sigma_i+r(\ell_i)p_{i-1}\le Cp_i^2$, and the fluctuation field obeys $|\phi_i'(x)|\le Cp_i^2$ there. (c) King's printed threshold ``$\dist(x,\partial T^{(k)})<p(L^k\eps)$'' in \cite[(5.43)--(5.45)]{K2} cannot be used: $p(\ell)=b_0(1+\ln\ell^{-1})^p$ with $2p\ge(d+1)\varrho$ (\cite[(3.54)]{Ba2}), so $p(\ell)/r(\ell)\to\infty$ and a contour of length up to $p(\ell)$ leaves $\Lambda_0^{(i)}$, where \cite[(5.41)]{K2} is imposed; the correct threshold is one corridor $r(\ell_i)$, with which the conclusions of (b) hold. (d) For $x$ at distance $>r(\ell_i)$ from $\DD$ the Dirichlet image effect on the background is $\le Ce^{-\delta_1r(\ell_i)}M_i\le p_i$ (\cref{lem:a01:FL,lem:a01:LOC}; $\varrho>1$), so \cite[(2.113)]{Ba2} holds there as printed. (e) Within distance $R_0$ of the face either a gradient violation or a face-family element is present on a shortest block path from a block of $R_s^{(i)}$ to the face layer, and at face blocks off the families the kept quartic is $\ge-C\lambda(\ell_i)p_i^8$ per block, an error $\le C\ell_i^\alpha$ per block; this is \cite[(5.46)--(5.49)]{K2} in event form.
\end{lemma}
\begin{proof}
(a) \Cref{lem:a01:FPC} with $\theta=\theta_2$ and $\sigma_i$ chosen so that $\frac12\theta_2\gamma_D\sigma_i^2/4=p_i^2/4$; the disjointness is the nesting of Ba{\l}aban's regions applied to a pointwise family (the face family is a condition on one block variable at one block of the index set, of the same type as $|\psi(y)-(Q\phi)(y)|>p$, so the rules (2.4)--(2.10) of \cite{Ba2} hold with it verbatim). (b) Triangle inequality along the path, the gradient clause $|D_{A_i}\phi_i|\le p_{i-1}$ on each of $\le r$ bonds and $|U|=1$; the fluctuation $\phi'=\phi_i-aL^{-2}C^{(i)}Q^*\phi_{i+1}$ has its second term bounded by $C\sup|\phi_{i+1}|$ on a neighbourhood (decay of $C^{(i)}$, \cref{lem:a01:A23}), which by (b) at scale $i+1$ is $\le Cp^2$. (c) Arithmetic: $p(\ell)/r(\ell)=(b_0/R)(1+\ln\ell^{-1})^{p-\varrho}$ with $p-\varrho=4$ at Ba{\l}aban's values $p=6$, $\varrho=2$. (d) \Cref{lem:a01:LOC} (sup form) times the block-sup bound $M_i$; $e^{-\delta R(1+\ln\ell^{-1})^2}M_i\le p(\ell)$ for $\ell\le\eta_1$. (e) Along a shortest cube-internal block path from $x\in R_s^{(i)}$ (where $|\phi_i(x)|\ge M_i$) to $\Lone^{(i)}$, of length $\le dR_0<s_b$ inside $\Lambda_7^{(i-1)\prime}$, either a bond carries a gradient violation, or the end block $y'$ has $|\phi_i(y')|\ge M_i-dR_0p_i\ge\frac12M_i\ge\sigma_i$ (for $\lambda(\ell_i)\le(2C_8)^{-4}$ and $dR_0\le\frac12\lambda(\ell_i)^{-1/4}$, conditions on $\eta_1$), a face-family element; the per-block error is $\lambda(\ell_i)(Cp_i^2)^4$ by (b).
\end{proof}

\begin{lemma}[Ba{\l}aban's per-scale lower bound at the weight]\label{lem:a01:329}
For a union $\Lambda$ of $k$-blocks with Neumann cuts and, where $\Lambda$ meets $\DD$, the Dirichlet condition of the cube, and $A$ regular on $\Lambda$ with $\theta_k\le\theta_*$,
\[
\ip\phi{\Delta^{(k)}(\Lambda,A)\phi}\ \ge\ \gamma_0^{w}\Big[\sum_{\text{block bonds }\langle y,y'\rangle\subset\Lambda}|U(A_k(\langle y,y'\rangle))\phi(y')-\phi(y)|^2+V_k\sum_y|\phi(y)|^2\Big]-C(\theta_k+e_kp_k)^2\sum_{y\in\Lambda}|\phi(y)|^2,
\]
with $\gamma_0^w>0$ depending on $(d,a,\Wt)$ only. On King's small-field set $|\phi|\le M_k'$ the last term is $\le Ce_k^2p_k^2M_k'^2|\Lambda|=O(\ell_k^{\kappa_0})|\Lambda|$; at $A=0$ it is absent. This is \cite[(3.29)]{Ba2} at the weight.
\end{lemma}
\begin{proof}
(1) Dirichlet faces only increase the form: the Dirichlet form on the live sites is the Neumann form plus the face terms $\vth^{d-2}|\phi(x)|^2\ge0$, so $\Delta^{(k)}_D\ge\Delta^{(k)}_N$ by the monotonicity of the minimiser identity. (2) $\ip\phi{\Delta\phi}=\min_\Phi[\ip\Phi{H\Phi}+a_k\norm{\phi-Q\Phi}^2]$. For a block bond $\langle y,y'\rangle$: $|\phi(y)-U\phi(y')|^2\le3(|\phi(y)-Q\Phi(y)|^2+|Q\Phi(y)-UQ\Phi(y')|^2+|\phi(y')-Q\Phi(y')|^2)$. On the two-block box $B(y)\cup B(y')$ (simply connected) gauge to the constant field and compare: the averages and covariant differences at $A$ and at the constant field differ by $\le C\theta_k\norm\Phi_{B\cup B'}$, and $\norm\Phi_B^2\le C(E_B(\Phi)+|\phi(y)|^2+|\phi(y)-Q\Phi(y)|^2)$ (block Poincar\'e with the weight); the block-bond transport $U(A_k)$ differs from the fine one by $Ce_kp_k$ (the photon block-versus-background step \cite[(A.2)--(A.6)]{K1}, which is $\kappa$-free and holds on the cube by \cref{lem:a01:PE,lem:a01:REG}). At $A=0$, $|\operatorname{avg}_wg-\operatorname{avg}_{w'}g|^2\le4\Wt C_P(d)E_{B\cup B'}(g)$ by Cauchy--Schwarz with the weights ($\sum\vth^dw^2=\Wt$ on a unit block) and the two-block Poincar\'e inequality. Mass: $V_k|\phi(y)|^2\le2V_k(\Wt\norm\Phi^2_{B(y)}+|\phi(y)-Q\Phi(y)|^2)$. Summing (each block term counted $\le2d$ times, each fine bond in $\le2$ boxes) and absorbing the $\theta_k$-parts gives the claim with $\gamma_0^w=[C(d)\max(4\Wt C_P,2\Wt,(6d+2)/a_k)]^{-1}(1-O(\theta_k))$.
\end{proof}

\begin{lemma}[the inherited face cap]\label{lem:a01:IC}
Let $k\ge1$ and fix a term of the step-$(k-1)$ large-field decomposition (a family of events with its bad union $S_{k-1}$). On the support of its characteristic functions (Ba{\l}aban's six families on $\Lambda_7^{(k-2)\prime}$ and the face family of scale $k-1$; at $k-1=0$ on the first live layer of fine sites), every face $k$-block $y\in\Lone^{(k)}$ with $B(y)\subset\Lambda_0^{(k-1)}$ satisfies, in unit-$k$ normalisation,
\[
|\phi_k(y)|\ \le\ L^{1/2}(\sigma_{k-1}+Lp_{k-1})\ \le\ C_{\mathrm{IC}}\,p_k,\qquad C_{\mathrm{IC}}\defeq L^{1/2}(1+\ln L)^p(C_8+L),
\]
free of $\eps$, $k$, $K$, $J$ and of $M$. By the nesting of the regions the cap holds at every face block of $\Lambda_2^{(k-1)\prime}$, in particular at every face block of the propagator domain $B^k(\Lambda_2^{(k-1)\prime})$, of the ring $\Lambda_6^{(k-1)\prime}\setminus\Lambda_7^{(k-1)\prime}$ and of $S_k$.
\end{lemma}
\begin{proof}
Every step-$(k-1)$ test located in $B(y)$ is imposed and passes: the face test at the $(k-1)$-blocks of $B(y)$ in $\Lone^{(k-1)}$ (their index set contains them since $\Lambda_0^{(k-1)}\subset\Lambda_7^{(k-2)\prime}$), the gradient test on the block bonds with both ends in $B(y)$ (a violating bond would put its endpoints in $S_{k-1}$), and the new-field test $|\phi_k(y)-(Q\phi_{k-1})(y)|\le p_{k-1}$ at $y$ ($y$ lies in the prime of $\Lambda_7^{(k-2)\prime}$ and $y\in P_s^{(k-1)}$ would give $B(y)\subset S_{k-1}$). Since $B(y)$ is a face block, some $(k-1)$-block $y''\subset B(y)$ is a face block with $|\phi_{k-1}(y'')|\le\sigma_{k-1}$; the gradient test along a block path inside $B(y)$ (length $\le L$) gives $|\phi_{k-1}|\le\sigma_{k-1}+Lp_{k-1}$ on all of $B(y)$, hence $|(Q\phi_{k-1})(y)|\le\sigma_{k-1}+Lp_{k-1}$ (Q5), and the new-field test gives $|\phi_k(y)|\le\sigma_{k-1}+(L+1)p_{k-1}$ in unit-$(k-1)$ normalisation, i.e.\ the claim after the conversion factor $L^{1/2}$ ($d=3$) and $p_{k-1}\le(1+\ln L)^pp_k$. At $k-1=0$ the step-$0$ face member is the cap on $F_1$ of \cref{def:a01:facefam}.
\end{proof}

Without the cap the fine-field class of \cite[(2.55)]{Ba2} fails at the collar cubes: for constant data of size $M_k$ the background's face gradient is of order $M_k$, not $p_k$. The cap is inherited from the previous step exactly as (2.55) is carried from step to step, as a characteristic function of the block field.
